\subsection{The Reduction Archive tab}
\label{sec:reduction_archive_tab}

All the products (final and intermediate) of all reduced configuration for all datasets are saved into
the {\tt EDPS\_data} directory. In addition, \edps\ offers the possibility to "archive" a desired reduction (e.g., the one
that gives the best results) by "exporting" only the final reduction products into another location.

Select the desired reduced configurations and press the "Archive button". This configuration disappears from the
Reduction Queue and appear in the tab "Reduction Archive".

The content of the "Reduction Archive", e.g. the list of archived reductions, can be inspected by pressing the
{\bf Reduction Archive} tab (Figure \ref{fig:archive_export_1}).



\begin{figure}[h!]
\centering
\includegraphics[width=0.9\textwidth]{figures/export1.jpg}
\caption{The reduction archive tab. This table contains all the different configurations of datasets that are declared "finished" and removed from the Reduction Queue. From this page, the user can export the most important files into a desired local directory.}
\label{fig:archive_export_1}
\end{figure}

Each configuration and its jobs can be inspected, unarchived and eventually deleted.
In order to copy the final products into a desired location, press the Export button. A window will appear as in Figure \ref{fig:archive_export_2}, 
allowing:

\begin{itemize}
 \item to select all or just the last configuration for the selected datasets  
 \item to specify a directory where to save the files. A check is done 
\end{itemize}
Press {\bf Export} to save the files.

The types of files to copy depend on each workflow. Typically, only most important files of the scientific tasks are saved. Other files can be found in the EDPS\_data directory.
The operation can take several minutes, depending on the sizes of the files.
The default structure is:

\verb|<dataset_name>/<reduction_time_stamp>/<category>.fits.|

\begin{figure}[h!]
\centering
\includegraphics[width=0.7\textwidth]{figures/export2.jpg}
\caption{The EXPORT dialogue window, where the user can decide which reduced configuration to save and where.}
\label{fig:archive_export_2}
\end{figure}
